% ============================================================================
\section{Ethics}\label{sec:ethics}
% ============================================================================

This study involved no human subjects. It measured what public storefront pages and the third-party code on them did; no person was observed, surveyed or contacted, and no personal data of real users were collected. The crawler used fresh, temporary browser profiles, never created accounts, logged in, added items to carts, submitted forms or made purchases, and interacted with consent banners only by clicking Decline or Reject, never Accept. Data flows were traced with a researcher-generated canary token. We used unmodified user agents, so our browser was visibly automated (Section~\ref{sec:app-validation}), and did not circumvent bot detection or access controls. The terms of service of individual stores were not reviewed.

\emph{\code{robots.txt}.} The Round 4 screening and crawls and ConsentCheck fetched each site's \code{robots.txt} before any other request and loaded a page only if it was allowed (Supplementary \ref*{S:sec:s-r4}). The main crawl did not consult \code{robots.txt} beforehand; a retrospective fetch found that 20 of the 1,427 stores that served one disallowed the exact landing URL we loaded, 16 only because of the campaign parameters that carried our canary (Supplementary \ref*{S:sec:s-sample}). These are the documented exceptions. Stores aimed at children or offering adult content were to be excluded, with one documented exception (Section~\ref{sec:method-sample}). After the crawl, a broader title screen found 2 adult-product stores that had been loaded, and a toy retailer and a second store that may be aimed at children were noticed; we excluded these 4 stores from all analyses for content sensitivity (Section~\ref{sec:method-sample}). Excluding in addition the 20 \code{robots.txt} cases changes no headline rate by more than 0.4 percentage points and leaves the main sample's blocked-on-paper count unchanged (Supplementary Table~\ref*{S:tab:s-ethics}).

\emph{Load on stores.} Including superseded, retried and re-check loads, each store received between 3 and 11 page loads (Supplementary \ref*{S:sec:s-sample}), at least 10.0~s apart. Most of the main crawl ran ten rather than the planned six concurrent browsers (Section~\ref{sec:method-deviations}), but each browser visited a different store, so the load on any single store did not change.

\emph{Records created by our visits.} Our visits may appear in merchants' analytics and ad-platform logs as sessions from a campaign named ``cmpstudy''; on creator links the token duplicated the creator's own campaign keys, which could affect attribution. Besides the refusals we clicked, some CMPs wrote consent states to Shopify on their own, in some cases because they detected an automated browser, and Shopify recorded a sale-of-data opt-out whenever our GPC header arrived (counts: Supplementary \ref*{S:sec:s-extended}). Future crawls should put a project-information URL in the token so that merchants can identify the study and opt out.

\emph{Reporting.} Results are reported in aggregate, and no store is named. The motivating store is called \storeC and is not named. CMP products are named where their documented capabilities determine the integration mode. The released data identify stores, and the sites of the Round 4 pilot, only by pseudonymous identifiers. This paper re-analyses pilot loads from our unpublished study of creator links, which reported the \storeC observation that motivated this one.

\emph{Disclosure.} The findings were disclosed on 2 October 2026 to Shopify, to the vendors of the three consent tools whose store-specific blocklists we evaluated (Pandectes, OneTrust and Usercentrics/Cookiebot) and to TikTok, with an invitation to correct Table~\ref{tab:cmp-modes}; the Pandectes disclosure included its merchant setting that grants consent to automated browsers. Pandectes reviewed the findings, confirmed our classification of its integration, and corrected two details of our description of its automated-browser setting (it writes no ``bot'' reason to Shopify, and the setting is off by default), which we adopted; OneTrust acknowledged receipt and asked for the store-level details, Usercentrics acknowledged receipt, and Shopify and TikTok had not replied by the time of submission.

\emph{Oversight.} No institutional review board or ethics committee determination was sought, because no human subjects were involved: the study observed no person and collected no personal data, only the behaviour of public web pages. The safeguards above (no accounts or personal data, public pages only, Decline clicks only, polite pacing, \code{robots.txt} with the documented exceptions, responsible disclosure, anonymized stores and a visibly automated browser) were applied throughout, including to the Round 4 crawls, which ran on 27 September 2026.

% ============================================================================
\section{Limitations}\label{sec:limits}
% ============================================================================

\emph{One vantage point, two days.} All loads came from one residential connection in California, on 26 September 2026 and, for Round 4, on 27 September 2026. Shopify's region defaults, and therefore the pre-choice results, differ for other US states and in particular for the EU and UK.

\emph{Visibly automated, headless browsers.} Every page could detect automation (Section~\ref{sec:app-validation}): where a consent tool grants consent to bots, pre-choice firing overstates what a person sees, and rules that hide the banner from bots lower the banner count and the detector's apparent recall, which the re-check, run with the same browser, cannot correct. Meta beacons are undercounted, and request-level sandbox shares may be higher than in a regular browser.

\emph{Declared policies.} Blocklists were readable for three CMP products. The primary blocked-on-paper estimate no longer rests on 14 stores but on 89, yet 80 of them use Pandectes, the expansion stores rank lower than the like-for-like reference, and the expansion screen stopped before its planned size (Section~\ref{sec:method-r4}). Integration modes are product-level labels, and no model can separate mode from product or default.

\emph{What we can observe.} We observed client-side requests only, so the data flows to vendors that we report are lower bounds; the store-level shares are not. A request that reached a vendor shows that data were sent, not how they were used. We never clicked Accept, the Decline arm had no repeat load, and its consent events were recorded on its last page only.

\emph{Labels, pilot and frame.} The validation labels have not yet been checked by an independent human rater (Section~\ref{sec:method-ai}). The pilot's sites were selected, not sampled, so its rates are not population estimates. The frame's recall is low in the top bucket, so the projections are conditional on the frame. Supplementary \ref*{S:sec:s-extended} gives the full list, including banner detection, the Decline procedure, timing and concurrency.

% ============================================================================
\section{Conclusion}\label{sec:conclusion}
The banner on \storeC did not govern the pixel that sent the video's code to TikTok; we measured how often and where such gaps open across 1,358 popular storefronts and 116 lower-ranked ones from California. Where a blocklist declared an ad vendor blocked and the consent tool had not itself consented, that vendor's pixel fired from Shopify's sandbox before any choice on 46 of 89 stores: sandboxed pixels followed Shopify's consent state, not the blocklist. Shopify's gate held back almost every pixel declaring a refused purpose, but its California default allowed tracking on most stores, automatic refusals came late, and pixels that started earlier or declared other purposes or none kept sending: after a working Decline, a new sandbox still sent ad requests on 23\% of stores, and GPC by design stopped only pixels declaring sale of data. A pilot on selected non-Shopify sites found ad requests after refusals too. The repair needs no new law: Shopify can hold its gate until the consent tool reports and stop pixels on withdrawal, consent tools can report every choice, and merchants can require consent and test their stores with ConsentCheck. Consent must be enforced by the component that starts the code.
